% Unit 036 terjemahan Bahasa Indonesia dari chapter3.tex baris 700--945.
% Sumber: Methods of Algebra, Volume 2, commit
% 9a5803ff77dd3257484cb177f851a73770a59dd3 (CC BY 4.0).
% stable-unit-id: o014.aljabr2.chapter3.double-complexes
% Cakupan: Bagian 3.5 lengkap.

% segment-id: o014.li2.chapter3.double-complexes.h001
\section{Bikompleks}\label{sec:double-cplx}
\hypertarget{o014.aljabr2.chapter3.double-complexes}{ }

% segment-id: o014.li2.chapter3.double-complexes.p001
Dalam bagian ini, $\mathcal{A}$ tetap merupakan kategori aditif. Secara kasar,
bikompleks dapat dibayangkan sebagai versi dua dimensi dari kompleks, dengan
diferensial dalam arah vertikal dan horizontal.

% segment-id: o014.li2.chapter3.double-complexes.d001
\begin{definisi}[Bikompleks]\label{def:double-cplx}
	\index{bikompleks (double complex)}
	\index[sym1]{dhoridvert@$\dhori$, $\dvert$}
	Bikompleks pada kategori aditif $\mathcal{A}$ terdiri atas keluarga objek
	$\left(X^{p,q}\right)_{(p,q)\in\Z^2}$ dalam $\mathcal{A}$, bersama dengan
	morfisme $\dhori^{p,q}:X^{p,q}\to X^{p+1,q}$ dan
	$\dvert^{p,q}:X^{p,q}\to X^{p,q+1}$, yang memenuhi
% segment-id: o014.li2.chapter3.double-complexes.q001
	\[ \dhori^{p+1,q} \dhori^{p,q} = 0, \quad \dvert^{p, q+1} \dvert^{p,q} = 0, \quad \dhori^{p, q+1} \dvert^{p,q} = \dvert^{p+1, q} \dhori^{p,q}. \]
	Data di atas biasanya disingkat sebagai
	$(X^{\bullet,\bullet},\dhori,\dvert)$, $X^{\bullet,\bullet}$, atau $X$.
\end{definisi}

% segment-id: o014.li2.chapter3.double-complexes.p002
Syarat-syarat pada $\dhori,\dvert$ dapat disingkat menjadi
% segment-id: o014.li2.chapter3.double-complexes.q002
\[ \dhori^2 = 0, \quad \dvert^2 = 0, \quad \dhori \dvert = \dvert \dhori. \]

% segment-id: o014.li2.chapter3.double-complexes.d002
\begin{definisi}\label{def:bicplx}
	Diberikan kategori aditif $\mathcal{A}$, suatu morfisme dari bikompleks $X$
	ke $Y$ adalah keluarga morfisme
% segment-id: o014.li2.chapter3.double-complexes.q003
	\[ f = \left( f^{p,q}: X^{p, q} \to Y^{p, q} \right)_{(p, q) \in \Z^2}, \]
	sedemikian sehingga, untuk semua $p,q$, berlaku
% segment-id: o014.li2.chapter3.double-complexes.q004
	\[ \dhori_Y^{p, q} f^{p, q} = f^{p+1, q} \dhori_X^{p, q}, \quad \dvert_Y^{p, q} f^{p, q} = f^{p, q+1} \dvert_X^{p, q}. \]
	Relasi-relasi di atas dapat disingkat sebagai
	$\dhori_Yf=f\dhori_X$ dan $\dvert_Yf=f\dvert_X$.
\end{definisi}

% segment-id: o014.li2.chapter3.double-complexes.p003
Seperti dalam Proposisi~\ref{prop:additive-cat-cplx}, definisi-definisi ini
menjadikan semua bikompleks pada $\mathcal{A}$ suatu kategori aditif
$\cate{C}^2(\mathcal{A})$.
\index[sym1]{C2A@$\cate{C}^2(\mathcal{A})$}

% segment-id: o014.li2.chapter3.double-complexes.p004
Secara visual, bikompleks $X$ digambarkan sebagai
% segment-id: o014.li2.chapter3.double-complexes.g001
\[\begin{tikzcd}
	& \vdots & \vdots & \\
	\cdots \arrow[r] & X^{p, q+1} \arrow[r] \arrow[u] & X^{p+1, q+1} \arrow[r] \arrow[u] & \cdots \\
	\cdots \arrow[r] & X^{p, q} \arrow[r, "{\dhori^{p,q}}"'] \arrow[u, "{\dvert^{p,q}}"] & X^{p+1, q} \arrow[r] \arrow[u] & \cdots \\
	& \vdots \arrow[u] & \vdots \arrow[u] &
\end{tikzcd}\]
Setiap baris $\left(X^{\bullet,q},d^{\bullet,q}\right)$ dan setiap kolom
$\left(X^{p,\bullet},d^{p,\bullet}\right)$ merupakan kompleks. Dengan
demikian, bikompleks dapat dipandang menurut kolom ataupun menurut baris.

% segment-id: o014.li2.chapter3.double-complexes.d003
\begin{definisi}\label{def:bicplx-F1}
	\index[sym1]{F1F2@$F_{\mathrm{I}}, F_{\mathrm{II}}$}
	Funktor aditif
	$F_{\mathrm{I}}:\cate{C}^2(\mathcal{A})\to
	\cate{C}(\cate{C}(\mathcal{A}))$ didefinisikan sebagai berikut: untuk
	bikompleks $X$, ambil $(F_{\mathrm{I}}X)^p=X^{p,\bullet}$ dan
	$d_{F_{\mathrm{I}}X}^p=\dhori^{p,\bullet}:(F_{\mathrm{I}}X)^p\to
	(F_{\mathrm{I}}X)^{p+1}$. Demikian pula, rumus
	$(F_{\mathrm{II}}X)^q=X^{\bullet,q}$ dan
	$d_{F_{\mathrm{II}}X}^q=\dvert^{\bullet,q}$ mendefinisikan funktor aditif
	$\cate{C}^2(\mathcal{A})\to\cate{C}(\cate{C}(\mathcal{A}))$. Di sini
	$p,q\in\Z$.
\end{definisi}

% segment-id: o014.li2.chapter3.double-complexes.p005
Definisi~\ref{def:bicplx} menyiratkan bahwa $F_{\mathrm{I}}$ dan
$F_{\mathrm{II}}$ keduanya merupakan isomorfisme antarkategori aditif. Hal
ini digambarkan sebagai berikut:
% segment-id: o014.li2.chapter3.double-complexes.g002
\begin{equation}\label{eqn:bicplx-diagram}\begin{tikzcd}[
		/tikz/execute at end picture={
			\node[rectangle, draw, fit=(NW) (SE), inner xsep=0.5em, inner ysep=0em] {};
		}]
		\vdots & |[alias=NW]| & {} & \\
		(F_{\mathrm{II}} X)^q \arrow[u, "{d^q_{F_{\mathrm{II}} X}}"] \arrow[phantom, r, ":=" description, sloped] & \cdots \arrow[r] & X^{p,q} \arrow[r, "{\dhori^{p,q}}"'] \arrow[u, "{\dvert^{p,q}}"] & \cdots \\
		\vdots \arrow[u] & & {} \arrow[u] & |[alias=SE]| \\
		& \cdots \arrow[r] & (F_{\mathrm{I}} X)^p \arrow[r, "{d^p_{F_{\mathrm{I}} X}}"] \arrow[phantom, u, ":=" description, sloped] & \cdots
\end{tikzcd}\end{equation}

% segment-id: o014.li2.chapter3.double-complexes.d004
\begin{definisi}[Kompleks total]\label{def:total-cplx}
	\index{kompleks total (total complex)}
	\index[sym1]{tot@$\tot_{\oplus}$, $\tot_{\Pi}$, $\tot$}
	Misalkan $\mathcal{A}$ mempunyai koproduk terhitung, yang dinotasikan
	dengan simbol jumlah langsung $\bigoplus$, dan misalkan $X$ objek
	$\cate{C}^2(\mathcal{A})$. Definisikan kompleks $\tot_{\oplus}(X)$ pada
	$\mathcal{A}$ sebagai berikut:
% segment-id: o014.li2.chapter3.double-complexes.l001
	\begin{compactitem}
% segment-id: o014.li2.chapter3.double-complexes.i001
		\item $\tot_{\oplus}(X)^n := \bigoplus_{p+q=n} X^{p, q}$,
% segment-id: o014.li2.chapter3.double-complexes.i002
		\item restriksi $d^n:\tot_{\oplus}(X)^n\to
		\tot_{\oplus}(X)^{n+1}$ pada suku $X^{p,q}$ sama dengan
		$\dhori^{p,q}+(-1)^p\dvert^{p,q}$.
	\end{compactitem}
\end{definisi}

% segment-id: o014.li2.chapter3.double-complexes.p006
Dengan menghilangkan superskrip, penulisan lengkap
$d^2:X^{p,q}\to X^{p+2,q}\oplus X^{p+1,q+1}\oplus X^{p,q+2}$ memberikan
% segment-id: o014.li2.chapter3.double-complexes.q005
\[ d^2 = \left( \dhori^2, \; (-1)^p \dhori \dvert + (-1)^{p+1} \dvert \dhori , \; \dvert^2 \right) = (0, 0, 0). \]

% segment-id: o014.li2.chapter3.double-complexes.p007
Jika koproduk dalam konstruksi funktor $\tot_{\oplus}$ diganti dengan produk
$\prod$, secara serupa diperoleh
$\tot_{\Pi}X\in\Obj(\cate{C}(\mathcal{A}))$, asalkan produk-produk yang
diperlukan ada. Definisi
$d^n:\tot_{\Pi}(X)^n\to\tot_{\Pi}(X)^{n+1}$ serupa dengan kasus
$\tot_{\oplus}$: kita mensyaratkan bahwa proyeksi $d^{n-1}$ ke $X^{p,q}$
sama dengan komposisi
% segment-id: o014.li2.chapter3.double-complexes.q006
\[ \tot_{\Pi}^{n-1}(X) \xrightarrow{\text{proyeksi}} X^{p-1, q} \oplus X^{p,q-1} \xrightarrow{(\dhori^{p-1, q}, (-1)^p \dvert^{p, q-1})} X^{p, q} \]
tersebut; dengan cara yang sama diperoleh $d^2=0$. Jika tidak menimbulkan
kerancuan, $\tot_{\oplus}X$ dan $\tot_{\Pi}X$ keduanya disebut
\emph{kompleks total} dari $X$.

% segment-id: o014.li2.chapter3.double-complexes.p008
Sifat berikut seharusnya sudah jelas.

% segment-id: o014.li2.chapter3.double-complexes.pr001
\begin{proposisi}
	Dengan syarat-syarat yang bersesuaian, $\tot_{\oplus}$ dan $\tot_{\Pi}$
	keduanya merupakan funktor aditif dari $\cate{C}^2(\mathcal{A})$ ke
	$\cate{C}(\mathcal{A})$.
\end{proposisi}

% segment-id: o014.li2.chapter3.double-complexes.p009
Dalam definisi kompleks total, superskrip $p$ dan $q$ sepintas tampak tidak
simetris. Untuk menghilangkan kesan ini, kita definisikan automorfisme aditif
$\mathrm{swap}$ pada $\cate{C}^2(\mathcal{A})$ sedemikian sehingga
$\mathrm{swap}(X)^{p,q}=X^{q,p}$ serta $\dhori$ dan $\dvert$ saling
dipertukarkan. Tanda $(-1)^{pq}$ yang akan menyertainya pada hakikatnya
merupakan mekanisme yang sama dengan aturan tanda Koszul dalam
\cite[Definisi--Teorema~7.4.4]{Li1}.

\index[sym1]{swap@$\mathrm{swap}$}
% segment-id: o014.li2.chapter3.double-complexes.pr002
\begin{proposisi}\label{prop:double-cplx-swap}
	Untuk setiap $(p,q)\in\Z^2$ dan objek $X$ dalam
	$\cate{C}^2(\mathcal{A})$, definisikan keluarga isomorfisme
% segment-id: o014.li2.chapter3.double-complexes.q007
	\[ r^{p, q}_X := (-1)^{pq} \identity_{X^{p,q}} : X^{p, q} \to \mathrm{swap}(X)^{q,p} . \]
% segment-id: o014.li2.chapter3.double-complexes.l002
	\begin{compactitem}
% segment-id: o014.li2.chapter3.double-complexes.i003
		\item Jika $\mathcal{A}$ mempunyai koproduk terhitung, maka
		$(r_X^{p,q})_{(p,q)\in\Z^2}$ menginduksi isomorfisme
		$r_X:\tot_{\oplus}(X)\rightiso\tot_{\oplus}(\mathrm{swap}(X))$.
% segment-id: o014.li2.chapter3.double-complexes.i004
		\item Jika $\mathcal{A}$ mempunyai produk terhitung, maka
		$(r_X^{p,q})_{(p,q)\in\Z^2}$ menginduksi isomorfisme
		$r_X:\tot_{\Pi}(X)\rightiso\tot_{\Pi}(\mathrm{swap}(X))$.
	\end{compactitem}
\end{proposisi}
% segment-id: o014.li2.chapter3.double-complexes.pf001
\begin{bukti}
	Untuk semua $(p,q)\in\Z^2$, diagram
% segment-id: o014.li2.chapter3.double-complexes.g003
	\[\begin{tikzcd}[row sep=large]
		X^{p,q} \arrow[d, "(-1)^{pq} \identity"'] \arrow[r, "{ (\dhori, (-1)^p \dvert) }" inner sep=0.8em] & X^{p+1, q} \oplus X^{p, q+1} \arrow[d, "{((-1)^{(p+1)q} \identity, (-1)^{p(q+1)} \identity) }"] \\
		X^{p,q} \arrow[r, "{((-1)^q \dhori, \dvert)}"' inner sep=0.8em] & X^{p+1, q} \oplus X^{p, q+1}
	\end{tikzcd}\]
	komutatif, sehingga $r_X$ memang merupakan morfisme kompleks.
\end{bukti}

% segment-id: o014.li2.chapter3.double-complexes.p010
Syarat adanya koproduk terhitung (atau produk terhitung) dalam definisi
$\tot_{\oplus}X$ (atau $\tot_{\Pi}X$) dapat diperlemah: jika untuk setiap
$n$ hanya ada paling banyak sejumlah hingga pasangan $(p,q)$ yang memenuhi
$p+q=n$ dan $X^{p,q}\neq0$, maka koproduk (atau produk) dalam definisi
kompleks total menjadi jumlah langsung hingga. Dalam hal ini, kedua kompleks
total dapat dinotasikan seragam sebagai $\tot(X)$. Pembahasan lebih lanjut
akan diberikan setelah Definisi~\sourcecrossref{def:Cf-Supp}{pada \S~3.10}.

% segment-id: o014.li2.chapter3.double-complexes.r001
\begin{catatan}\label{rem:k-fold-cplx}
	Dengan gagasan yang sama, untuk setiap $k\in\Z_{\geq1}$ dapat didefinisikan
	kompleks lipat-$k$ sebagai data
% segment-id: o014.li2.chapter3.double-complexes.q008
	\[ \left( X^{p_1, \ldots, p_k} \in \Obj(\mathcal{A}) \right)_{p_1, \ldots, p_k \in \Z}, \quad {}^1 d, \ldots, {}^k d, \]
	dengan ${}^i d^{p_1,\ldots,p_k}:X^{p_1,\ldots,p_k}\to
	X^{p_1,\ldots,p_i+1,\ldots,p_k}$ dan
% segment-id: o014.li2.chapter3.double-complexes.q009
	\[ {}^i d^2 = 0, \quad {}^i d {}^j d = {}^j d {}^i d \quad \text{(superskrip dihilangkan)}. \]
	Semua kompleks lipat-$k$ membentuk kategori aditif
	$\cate{C}^k(\mathcal{A})$. Dengan mengikuti
	Definisi~\ref{def:bicplx-F1}, dapat pula didefinisikan keluarga funktor
	aditif $F_i:\cate{C}^k(\mathcal{A})\to
	\cate{C}(\cate{C}^{k-1}(\mathcal{A}))$ sedemikian sehingga setiap $F_i$
	merupakan ekuivalensi antarkategori ($k\geq2$ dan $i=1,\ldots,k$).
	
	Dalam keadaan ini, jika $\mathcal{A}$ mempunyai koproduk atau produk
	terhitung, funktor kompleks total dapat didefinisikan dengan cara serupa:
% segment-id: o014.li2.chapter3.double-complexes.q010
	\[ \tot_{\oplus} \; \text{atau} \; \tot_{\Pi}: \cate{C}^k(\mathcal{A}) \to \cate{C}(\mathcal{A}). \]
	Rinciannya sama dengan kasus bikompleks dan tidak perlu diuraikan lagi.
\end{catatan}

% segment-id: o014.li2.chapter3.double-complexes.p011
Definisikan isomorfisme dari $\cate{C}^2(\mathcal{A})$ ke dirinya sendiri
$[m]_{\mathrm{I}}:=F_{\mathrm{I}}^{-1}\circ[m]\circ F_{\mathrm{I}}$ dan
$[m]_{\mathrm{II}}:=F_{\mathrm{II}}^{-1}\circ[m]\circ F_{\mathrm{II}}$,
yang masing-masing bersesuaian dengan pergeseran horizontal dan vertikal
bikompleks. Untuk semua $a,b\in\Z$ berlaku
$[a]_{\mathrm{I}}[b]_{\mathrm{II}}=[b]_{\mathrm{II}}[a]_{\mathrm{I}}$.
Sekarang kita kaji pengaruh funktor-funktor ini terhadap kompleks total serta
hubungannya dengan automorfisme $\mathrm{swap}$.

% segment-id: o014.li2.chapter3.double-complexes.n001
% Reference: KS06, p.290
% segment-id: o014.li2.chapter3.double-complexes.pr003
\begin{proposisi}\label{prop:tot-shift}
	Untuk semua $X\in\Obj(\cate{C}^2(\mathcal{A}))$ dan
	$(p,q)\in\Z^2$, notasikan pembenaman standar
	$X^{p,q}\hookrightarrow\mathrm{tot}_{\oplus}(X)^{p+q}$ dengan
	$\iota_{p,q}$. Untuk setiap $n\in\Z$, definisikan
% segment-id: o014.li2.chapter3.double-complexes.l003
	\begin{itemize}
% segment-id: o014.li2.chapter3.double-complexes.i005
		\item $\theta_X^n:\tot_{\oplus}\left(X[1]_{\mathrm{I}}\right)^n\to
		\tot_{\oplus}(X)[1]^n$ sedemikian sehingga restriksinya pada
		$X[1]_{\mathrm{I}}^{p,q}=X^{p+1,q}$ sama dengan $\iota_{p+1,q}$;
% segment-id: o014.li2.chapter3.double-complexes.i006
		\item $(\theta'_X)^n:\tot_{\oplus}\left(X[1]_{\mathrm{II}}\right)^n
		\to\tot_{\oplus}(X)[1]^n$ sedemikian sehingga restriksinya pada
		$X[1]_{\mathrm{II}}^{p,q}=X^{p,q+1}$ sama dengan
		$(-1)^p\iota_{p,q+1}$.
	\end{itemize}
	Keduanya memberikan isomorfisme kanonik dalam $\cate{C}(\mathcal{A})$:
% segment-id: o014.li2.chapter3.double-complexes.q011
	\begin{align*}
		\theta_X & = (\theta_X^n)_n: \tot_{\oplus} \left( X[1]_{\mathrm{I}} \right) \rightiso \tot_{\oplus} \left(X\right)[1], \\
		\theta'_X & = ((\theta'_X)^n)_n: \tot_{\oplus} \left( X[1]_{\mathrm{II}} \right) \rightiso \tot_{\oplus} \left(X\right)[1],
	\end{align*}
	beserta diagram antikomutatif (yakni, kedua komposisinya berselisih tanda
	negatif)
% segment-id: o014.li2.chapter3.double-complexes.g004
	\[\begin{tikzcd}
		\tot_{\oplus} \left( X[1]_{\mathrm{I}}[1]_{\mathrm{II}}\right) \arrow[r, "{\theta_{X[1]_{\mathrm{II}}}}" inner sep=0.6em] \arrow[d, "{\theta'_{X[1]_{\mathrm{I}}}}"'] & \tot_{\oplus}\left(X[1]_{\mathrm{II}}\right)[1] \arrow[d, "{\theta'_X [1]}"] \\
		\tot_{\oplus} \left( X[1]_{\mathrm{I}}\right)[1] \arrow[r, "{\theta_X [1]}"' inner sep=0.6em] & \tot_{\oplus}(X)[2]
	\end{tikzcd}\]
	dan diagram komutatif
% segment-id: o014.li2.chapter3.double-complexes.g005
	\[\begin{tikzcd}[column sep=large]
		\tot_{\oplus} \left( X[1]_{\mathrm{I}}[1]_{\mathrm{II}} \right) \arrow[r, "\text{$\theta$ lalu $\theta'$}"] \arrow[d, "{r_{X[1]_{\mathrm{I}}[1]_{\mathrm{II}}}}"'] & \tot_{\oplus}(X)[2] \arrow[d, "{r_X[2]}"] \\
		\tot_{\oplus}\left( \mathrm{swap}(X)[1]_{\mathrm{I}}[1]_{\mathrm{II}} \right) \arrow[r, "\text{$\theta'$ lalu $\theta$}"'] & \tot_{\oplus}\left(\mathrm{swap}(X)\right)[2],
	\end{tikzcd}\]
	dengan $r$ sebagaimana didefinisikan dalam
	Proposisi~\ref{prop:double-cplx-swap}. Jika $\tot_{\oplus}$ diganti dengan
	$\tot_{\Pi}$, tetap diperoleh $\theta_X,\theta'_X$ yang sama beserta diagram
	antikomutatif yang bersesuaian.
\end{proposisi}
% segment-id: o014.li2.chapter3.double-complexes.pf002
\begin{bukti}
	Verifikasi langsung berdasarkan definisi kompleks total
	(Definisi~\ref{def:total-cplx}) dan definisi funktor pergeseran
	(Definisi~\ref{def:translation-functor}) menunjukkan bahwa $\theta_X$ dan
	$\theta'_X$ merupakan morfisme kompleks; rinciannya panjang tetapi tidak
	sulit. Diagram antikomutatif tereduksi pada pengamatan bahwa diagram
% segment-id: o014.li2.chapter3.double-complexes.g006
	\[\begin{tikzcd}
		X^{p+1, q+1} \arrow[d, "{(-1)^p \identity}"'] \arrow[r, "\identity"] & X^{p+1, q+1} \arrow[d, "{(-1)^{p+1} \identity}"] \\
		X^{p+1, q+1} \arrow[r, "\identity"'] & X^{p+1, q+1}
	\end{tikzcd}\]
	antikomutatif.
	
	Sekarang tinjau pernyataan tentang diagram komutatif. Tetapkan
	$(p,q)\in\Z^2$ dan bandingkan restriksi kedua komposisi pada
	$(X[1]_{\mathrm{I}}[1]_{\mathrm{II}})^{p,q}$. Telah dijelaskan bahwa
	$\theta$ lalu $\theta'$ (untuk $X$) dan $\theta'$ lalu $\theta$ (untuk
	$\mathrm{swap}(X)$) masing-masing memberikan
% segment-id: o014.li2.chapter3.double-complexes.q012
	\[ (-1)^{p+1}, \; (-1)^q \; \in \Aut(X^{p+1, q+1}) \quad \text{($\identity$ dihilangkan)}. \]
	Di sisi lain, restriksi $r_{X[1]_{\mathrm{I}}[1]_{\mathrm{II}}}$ dan
	$r_X[2]$ pada $X^{p+1,q+1}$ masing-masing adalah $(-1)^{pq}$ dan
	$(-1)^{(p+1)(q+1)}$. Namun,
	$(p+1)(q+1)-pq\equiv(p+1)-q\pmod{2}$. Terbukti.
\end{bukti}

% segment-id: o014.li2.chapter3.double-complexes.p012
Sebagai perumuman sederhana dari diagram antikomutatif di atas, pembaca
diminta membuktikan bahwa, untuk semua $a,b\in\Z$, kedua morfisme
% segment-id: o014.li2.chapter3.double-complexes.q013
\[ \tot_{\oplus}(X[a]_{\mathrm{I}}[b]_{\mathrm{II}}) \rightrightarrows \tot_{\oplus}(X)[a+b] \]
yang didefinisikan dengan kedua urutan ($\theta$ lalu $\theta'$ dan $\theta'$
lalu $\theta$) berselisih tanda $(-1)^{ab}$; di sisi lain, diagram komutatif
di atas tetap berlaku untuk $X[a]_{\mathrm{I}}[b]_{\mathrm{II}}$. Pernyataan
yang sama berlaku jika $\tot_{\oplus}$ diganti dengan $\tot_{\Pi}$.

% segment-id: o014.li2.chapter3.double-complexes.p013
Salah satu penggunaan utama bikompleks ialah dalam kajian bifunktor. Apakah
yang dimaksud dengan bifunktor?

% segment-id: o014.li2.chapter3.double-complexes.c001
\begin{konvensi}\label{con:bifunctor}
	\index{bifunktor (bifunctor)}
	Funktor berbentuk $F:\mathcal{A}_1\times\mathcal{A}_2\to\mathcal{B}$
	disebut \emph{bifunktor}. Setelah objek
	$X_i\in\Obj(\mathcal{A}_i)$ ditetapkan, diperoleh funktor satu variabel
	$F(X_1,\cdot):\mathcal{A}_2\to\mathcal{B}$ dan
	$F(\cdot,X_2):\mathcal{A}_1\to\mathcal{B}$. Dengan demikian, dapat dibahas
	keaditifan, keeksakan, dan berbagai konsep lain bagi setiap variabel $F$.
\end{konvensi}

% segment-id: o014.li2.chapter3.double-complexes.dp001
\begin{definisiproposisi}\label{def:bifunctor-cplx}
	\index[sym1]{C2F@$\cate{C}^2 F, \cate{C}_{\oplus} F, \cate{C}_{\Pi} F$}
	Misalkan $\mathcal{A}_1,\mathcal{A}_2,\mathcal{B}$ kategori-kategori
	aditif dan $F:\mathcal{A}_1\times\mathcal{A}_2\to\mathcal{B}$ suatu
	bifunktor yang aditif dalam setiap variabel. Terdapat funktor yang
	bersesuaian
% segment-id: o014.li2.chapter3.double-complexes.g007
	\[\begin{tikzcd}[row sep=tiny, column sep=small]
		\cate{C}^2 F: \cate{C}(\mathcal{A}_1) \times \cate{C}(\mathcal{A}_2) \arrow[r] & \cate{C}^2(\mathcal{B}) \\
		(X, Y) \arrow[mapsto, r] & (F(X, Y))^{p, q} := F(X^p, Y^q), \\
		& \dhori^{p, q} = F\left( d_X^p, \identity \right), \; \dvert^{p, q} = F\left( \identity, d_Y^q\right) .
	\end{tikzcd}\]

	Oleh karena itu, jika $\mathcal{B}$ mempunyai koproduk terhitung atau
	produk terhitung, masing-masing dapat didefinisikan
% segment-id: o014.li2.chapter3.double-complexes.q014
	\begin{gather*}
		\cate{C}_{\oplus} F := \tot_{\oplus} \circ \cate{C}^2 F: \cate{C}(\mathcal{A}_1) \times \cate{C}(\mathcal{A}_2) \to \cate{C}(\mathcal{B}), \\
		\cate{C}_{\Pi} F := \tot_{\Pi} \circ \cate{C}^2 F: \cate{C}(\mathcal{A}_1) \times \cate{C}(\mathcal{A}_2) \to \cate{C}(\mathcal{B}).
	\end{gather*}
	Jika $\tot_{\oplus}$ dan $\tot_{\Pi}$ yang dibahas hanya melibatkan
	jumlah langsung hingga, maka $\cate{C}_{\oplus}F=\cate{C}_{\Pi}F$.%
	\footnote{Catatan penerjemah (O014-C038): sumber menghilangkan $F$ dari
	kesetaraan ini dan dari suku paling kanan pada dua isomorfisme di bawah.
	Karena funktor yang didefinisikan ialah $\cate{C}_{\oplus}F$ dan
	$\cate{C}_{\Pi}F$, target memulihkan $F$ pada ketiga tempat tersebut.}
\end{definisiproposisi}
% segment-id: o014.li2.chapter3.double-complexes.pf003
\begin{bukti}
	Syarat bikompleks $\dhori\dvert=\dvert\dhori$ langsung mengikuti definisi
	bifunktor; selebihnya jelas.
\end{bukti}

% segment-id: o014.li2.chapter3.double-complexes.p014
Kategori bikompleks juga mempunyai konsep homotopi. Misalkan
$f,g:X\to Y$ sepasang morfisme dalam $\cate{C}^2(\mathcal{A})$. Homotopi
dari $f$ ke $g$ adalah dua keluarga morfisme yang memenuhi syarat-syarat
berikut. \index{homotopi!versi bikompleks}
% segment-id: o014.li2.chapter3.double-complexes.q015
\begin{equation*}\begin{gathered}
% segment-id: o014.li2.chapter3.double-complexes.g008
\begin{tikzcd}
	Y^{p-1, q} & X^{p,q} \arrow[l, "{h^{p,q}}"'] \arrow[d, "{k^{p, q}}"] \\
	& Y^{p, q-1}
\end{tikzcd} \qquad (p,q) \in \Z^2 , \\
	\dvert^{p-1, q} h^{p, q} = h^{p, q+1} \dvert^{p, q}, \qquad \dhori^{p, q-1} k^{p,q} = k^{p+1, q} \dhori^{p,q}, \\
	g^{p,q} - f^{p, q} = \dhori^{p-1, q} h^{p,q} + h^{p+1, q} \dhori^{p,q} + \dvert^{p, q-1} k^{p,q} + k^{p, q+1} \dvert^{p,q}.
\end{gathered}\end{equation*}
Hal ini beralasan: dari $\dvert h=h\dvert$, $\dhori k=k\dhori$, dan definisi
bikompleks, mudah diperiksa bahwa
$\dhori h+h\dhori+\dvert k+k\dvert$ selalu memberikan morfisme dalam
$\cate{C}^2(\mathcal{A})$.

% segment-id: o014.li2.chapter3.double-complexes.p015
Homotopi bikompleks tercermin pada kompleks total. Lebih tepatnya, setelah
$(h^{p,q},k^{p,q})_{p,q\in\Z}$ diberikan, dapat didefinisikan
$\hat{h}\in\Hom^{-1}\left(\tot_{\oplus}(X),\tot_{\oplus}(Y)\right)$
sedemikian sehingga restriksi $\hat{h}$ pada suku jumlah langsung $X^{p,q}$
sama dengan
% segment-id: o014.li2.chapter3.double-complexes.q016
\[ \left( h^{p,q}, (-1)^p k^{p,q} \right): X^{p,q} \to Y^{p-1, q} \oplus Y^{p, q-1}, \]
yang menjadikan
$\tot_{\oplus}(g)-\tot_{\oplus}(f)=d^{-1}\hat{h}$; kasus $\tot_{\Pi}$
ditangani dengan cara yang sama.

% segment-id: o014.li2.chapter3.double-complexes.p016
Dengan demikian, relasi homotopi bikompleks mendefinisikan
$\cate{K}^2(\mathcal{A})$ beserta funktor
$\cate{C}^2(\mathcal{A})\to\cate{K}^2(\mathcal{A})$. Jika koproduk terhitung
(atau produk terhitung) ada, funktor $\tot_{\oplus}$ (atau $\tot_{\Pi}$)
turun menjadi funktor
$\cate{K}^2(\mathcal{A})\to\cate{K}(\mathcal{A})$.
\index[sym1]{K2A@$\cate{K}^2(\mathcal{A})$}

% segment-id: o014.li2.chapter3.double-complexes.p017
Kembali ke bifunktor $F:\mathcal{A}_1\times\mathcal{A}_2\to\mathcal{B}$
di antara kategori-kategori aditif, dan andaikan $F$ aditif dalam setiap
variabel. Berikut adalah versi bifunktor dari \eqref{eqn:KF}.

% segment-id: o014.li2.chapter3.double-complexes.pr004
\begin{proposisi}\label{prop:bifunctor-cplx-homotopy}
	\index[sym1]{K2F@$\cate{K}^2 F, \cate{K}_{\oplus} F, \cate{K}_{\Pi} F$}
	Ambil bifunktor $F$ seperti di atas. Maka $\cate{C}^2F$ berfaktor melalui
	$\cate{K}^2F:\cate{K}(\mathcal{A}_1)\times\cate{K}(\mathcal{A}_2)\to
	\cate{K}^2(\mathcal{B})$.
	
	Demikian pula, $\cate{C}_{\oplus}F$ atau $\cate{C}_{\Pi}F$ berfaktor
	melalui $\cate{K}(\mathcal{A}_1)\times\cate{K}(\mathcal{A}_2)\to
	\cate{K}(\mathcal{B})$, dan masing-masing dinotasikan dengan
	$\cate{K}_{\oplus}F$ atau $\cate{K}_{\Pi}F$, asalkan funktor
	$\cate{C}_{\oplus}F$ atau $\cate{C}_{\Pi}F$ terdefinisi.
\end{proposisi}
% segment-id: o014.li2.chapter3.double-complexes.pf004
\begin{bukti}
	Ambil kasus $\cate{C}_{\oplus}F$ sebagai contoh. Yang perlu ditunjukkan
	ialah bahwa, untuk morfisme homotopik nol
	$f_1=d^{-1}g:X_1\to Y_1$ dalam $\cate{C}(\mathcal{A}_1)$ dan sebarang
	morfisme $f_2:X_2\to Y_2$ dalam $\cate{C}(\mathcal{A}_2)$, morfisme yang
	bersesuaian $(\cate{C}^2F)(f_1,f_2):\cate{C}^2F(X_1,X_2)\to
	\cate{C}^2F(Y_1,Y_2)$ homotopik nol. Pilihan yang jelas ialah
	$h^{p,q}:=F(g^p,f_2^q)$ dan $k^{p,q}:=0$. Karena $f_2$ merupakan morfisme,
	$h$ memang berkomutasi dengan $\dvert=F(\identity,d)$. Variabel kedua
	ditangani dengan cara serupa.
\end{bukti}

% segment-id: o014.li2.chapter3.double-complexes.p018
Proposisi~\ref{prop:tot-shift} memberikan isomorfisme kanonik
% segment-id: o014.li2.chapter3.double-complexes.q017
\begin{gather*}
	\cate{C}_{\oplus} F(X[1], Y) \simeq \cate{C}_{\oplus}F(X, Y)[1] \simeq \cate{C}_{\oplus}F(X, Y[1]), \\
	\cate{K}_{\oplus} F(X[1], Y) \simeq \cate{K}_{\oplus}F(X, Y)[1] \simeq \cate{K}_{\oplus}F(X, Y[1]).
\end{gather*}
Pernyataan yang sama berlaku jika $\oplus$ diganti dengan $\Pi$.

% segment-id: o014.li2.chapter3.double-complexes.e001
\begin{contoh}[Bikompleks $\Hom$]\label{eg:Hom-bicplx}
	\index{bikompleks Hom@bikompleks $\Hom$ ($\Hom$ double complex)}
	\index[sym1]{Hombb@$\Hom^{\bullet, \bullet}$}
	Tinjau bifunktor
	$\Hom(\cdot,\cdot):\mathcal{A}^{\opp}\times\mathcal{A}\to\cate{Ab}$,
	yang memetakan pasangan objek $(S,T)$ ke $\Hom_{\mathcal{A}}(S,T)$.
	Berikut ini akan ditunjukkan bahwa kompleks $\Hom$, yaitu
	$\Hom^\bullet(X,Y)$, isomorfik secara kanonik dengan nilai pada objek
	$(X,Y)$ dari funktor komposit
% segment-id: o014.li2.chapter3.double-complexes.q018
	\[ \cate{C}(\mathcal{A})^{\opp} \times \cate{C}(\mathcal{A}) \xrightarrow{(\sigma^{-1}, \identity)} \cate{C}(\mathcal{A}^{\opp}) \times \cate{C}(\mathcal{A}) \xrightarrow{\cate{C}_{\Pi} \Hom(\cdot, \cdot)} \cate{C}(\cate{Ab}) \]
	dengan $\sigma$ sebagaimana dalam
	Definisi--Proposisi~\ref{def:sigma}.

	Untuk itu, pertama-tama tinjau bifunktor
	$F:\mathcal{A}\times\mathcal{A}^{\opp}\to\cate{Ab}$ yang didefinisikan
	dengan $F(T,S)=\Hom_{\mathcal{A}}(S,T)$, serta funktor komposit
% segment-id: o014.li2.chapter3.double-complexes.q019
	\[ \cate{C}(\mathcal{A}) \times \cate{C}(\mathcal{A})^{\opp} \xrightarrow{(\identity, \sigma^{-1})} \cate{C}(\mathcal{A}) \times \cate{C}(\mathcal{A}^{\opp}) \xrightarrow{\cate{C}^2 F} \cate{C}^2(\cate{Ab}). \]
	Misalkan $X,Y\in\Obj(\cate{C}(\mathcal{A}))$. Citra objek $(Y,X)$ oleh
	funktor komposit di atas disebut \emph{bikompleks $\Hom$} dan dinotasikan
	dengan $\Hom^{\bullet,\bullet}(X,Y)$. Berdasarkan diagram komutatif yang
	jelas
% segment-id: o014.li2.chapter3.double-complexes.g009
	\[\begin{tikzcd}[column sep=large]
		\cate{C}(\mathcal{A})^{\opp} \times \cate{C}(\mathcal{A}) \arrow[d, "\simeq"', "\text{pertukaran}"] \arrow[r, "{(\sigma^{-1}, \identity)}" inner sep=0.6em] & \cate{C}(\mathcal{A}^{\opp}) \times \cate{C}(\mathcal{A}) \arrow[r, "{\cate{C}^2 \Hom(\cdot, \cdot)}" inner sep=0.6em] \arrow[d, "\simeq"', "\text{pertukaran}"] & \cate{C}^2(\cate{Ab}) \arrow[d, "\mathrm{swap}"] \\
		\cate{C}(\mathcal{A}) \times \cate{C}(\mathcal{A})^{\opp} \arrow[r, "{(\identity, \sigma^{-1})}"' inner sep=0.6em] & \cate{C}(\mathcal{A}) \times \cate{C}(\mathcal{A}^{\opp}) \arrow[r, "{\cate{C}^2 F}"' inner sep=0.6em] & \cate{C}^2(\cate{Ab})
	\end{tikzcd}\]
	dan Proposisi~\ref{prop:double-cplx-swap}, persoalan semula tereduksi pada
	pembuktian adanya isomorfisme kanonik
	\[
		\tot_{\Pi}\Hom^{\bullet,\bullet}(X,Y)
		\rightiso \Hom^\bullet(X,Y).
	\]
	
	Dengan mencermati definisi, diperoleh
% segment-id: o014.li2.chapter3.double-complexes.q020
	\begin{equation*}\begin{aligned}
		\Hom^{p, q}(X, Y) & = \Hom_{\mathcal{A}}\left( X^{-q}, Y^p \right), \\
		\dhori^{p, q} & = (d_Y^p)_* : \Hom_{\mathcal{A}}\left( X^{-q}, Y^p \right) \to \Hom_{\mathcal{A}}\left( X^{-q}, Y^{p+1} \right), \\
		\dvert^{p, q} & = (-1)^q (d_X^{-q-1})^* : \Hom_{\mathcal{A}}\left( X^{-q}, Y^p \right) \to \Hom_{\mathcal{A}}\left( X^{-q-1}, Y^p \right).
	\end{aligned}\end{equation*}%
	\footnote{Catatan penerjemah (O014-C039): sumber memakai koefisien
	$(-1)^{q+1}$, yang sesuai dengan penggunaan ulang rumus $\sigma$ tetapi
	bukan dengan invers ketat $\sigma^{-1}$ yang diperoleh dari koreksi tanda
	pada bagian sebelumnya. Target memakai koefisien invers
	$(-1)^q$; akibatnya, totalisasi diidentifikasi dengan kompleks $\Hom$
	melalui tanda komponen $(-1)^q$, bukan melalui kesamaan harfiah.}
	Sekarang tentukan $\tot_{\Pi}\Hom^{\bullet,\bullet}(X,Y)$. Pertama,
% segment-id: o014.li2.chapter3.double-complexes.q021
	\[ \left(\tot_{\Pi} \Hom^{\bullet, \bullet}(X, Y)\right)^n = \prod_{p+q=n} \Hom_{\mathcal{A}}\left( X^{-q}, Y^p \right) = \prod_{k \in \Z} \Hom\left( X^k, Y^{k+n} \right) \]
	(dengan substitusi $k=-q$), yang tidak lain ialah $\Hom^n(X,Y)$.
	Selanjutnya, deskripsikan
	$d_{\tot_{\Pi}\Hom^{\bullet,\bullet}(X,Y)}^n$: koordinat $(p,q)$-nya
	dalam $\Hom^{n+1}(X,Y)$, dengan $p+q=n+1$, berasal dari
% segment-id: o014.li2.chapter3.double-complexes.g010
	\[\begin{tikzcd}[column sep=large]
		\Hom^{p-1, q}(X, Y) \arrow[r, "{\dhori^{p-1, q}}"] & \Hom^{p, q}(X, Y) \\
		& \Hom^{p, q-1}(X, Y) \arrow[u, "{(-1)^p \dvert^{p, q-1}}"'] .
	\end{tikzcd}\]
	Panah horizontal adalah $(d_Y^{p-1})_*$, sedangkan panah vertikal adalah
	$(-1)^{p+q-1}(d_X^{-q})^*=(-1)^n(d_X^{-q})^*$. Isomorfisme yang
	mengalikan faktor $\Hom^{p,q}(X,Y)$ dengan $(-1)^q$ mengubah tanda suku
	vertikal. Dibandingkan dengan Definisi~\ref{def:Hom-cplx}, hal ini
	memberikan isomorfisme kanonik
	$\tot_{\Pi}\Hom^{\bullet,\bullet}(X,Y)\rightiso\Hom^\bullet(X,Y)$.
	Terbukti.

	Terakhir, jika $\mathcal{A}$ kategori $\Bbbk$-linear, maka
	$\Hom^{\bullet,\bullet}(X,Y)$ meningkat menjadi funktor bernilai dalam
	$\cate{C}^2(\Bbbk\dcate{Mod})$.
\end{contoh}
